\documentclass[10pt]{article}
\usepackage{amsmath,amssymb,amsthm}
\usepackage[margin=1.05in]{geometry}

\newtheorem{theorem}{Theorem}[section]
\newtheorem{lemma}[theorem]{Lemma}
\newtheorem{proposition}[theorem]{Proposition}
\newtheorem{corollary}[theorem]{Corollary}
\newtheorem{conjecture}[theorem]{Conjecture}
\theoremstyle{definition}
\newtheorem{definition}[theorem]{Definition}
\theoremstyle{remark}
\newtheorem{remark}[theorem]{Remark}
\newtheorem{observation}[theorem]{Observation}

\newcommand{\Z}{\mathbb{Z}}
\newcommand{\N}{\mathbb{N}}
\newcommand{\floor}[1]{\left\lfloor #1 \right\rfloor}
\newcommand{\ceil}[1]{\left\lceil #1 \right\rceil}
\newcommand{\lb}{\log_2 3}

\title{The Terras coefficient-stopping conjecture\\ and the Pythagorean comma}
\author{Martien de Jong\\[2pt]
\normalsize with computational collaboration by Jengo (AI)}
\date{July 2026}

\begin{document}
\maketitle

\begin{abstract}
Terras (1976) attached to each positive integer $n$ two quantities: the stopping
time $\sigma(n)$, the first iterate of the accelerated Collatz map that falls
below $n$, and the coefficient stopping time $\tau(n)$, the first time the
accumulated multiplicative coefficient $3^{u}/2^{t}$ falls below $1$. Always
$\tau(n) \le \sigma(n)$; the \emph{coefficient stopping time conjecture} (CST)
asserts $\tau(n) = \sigma(n)$ for all $n > 1$. We verify CST exhaustively over
all $81{,}119$ stopping classes of depth $\le 24$ (covering every $n$ with
$\tau(n) \le 24$; the sole sub-threshold class member is $n = 1$, the trivial
cycle, an exact equality case) and directly for all odd $1 < n \le 10^{6}$
(zero violations; largest $\tau$ observed $176$). A violation in the stopping
class with $u$ odd steps among $j$ total steps and additive constant $b$
requires $n \le b/(2^{j}-3^{u})$. We prove the unconditional bound
$b/2^{j} < u/3$, from which the verified basis yields
$\tau(n)=\sigma(n)$ \emph{for all $n>1$ with $\tau(n) \le 4700$}, and, beyond
that depth, confines every possible failure of CST to an explicit thin set of
classes: those whose $(u,j)$ is a best rational approximation to $\lb$ from
above. The smallest safety margins concentrate at $(u,j)=(5,8)$, the
Pythagorean interval $2^{8}/3^{5} = 256/243$, the first nontrivial upper
convergent of $\lb$; the next danger zones are $(41,65)$, $(306,485)$, and
$(15601,24727)$. The maximal constant $B(t) = \max b/2^{j}$ over stopping
classes of depth $t$ is measured to grow like $1 + t/8$, one unit per comma
block, with the extremal classes given by repeated comma words; we prove the
weaker linear bound $B(t) < t/(3\lb)$, which suffices for all reductions used.
Consequently CST is reduced to effective rational approximation of $\lb$ ---
the same transcendence wall (Baker, Rhin) that governs the exclusion of
nontrivial cycles. One Diophantine question controls both classical
subproblems.
\end{abstract}

\section{Introduction}

Let $T$ denote the accelerated Collatz (``Terras'') map on the positive
integers,
\[
T(n) \;=\;
\begin{cases}
 n/2, & n \equiv 0 \pmod 2,\\[2pt]
 (3n+1)/2, & n \equiv 1 \pmod 2.
\end{cases}
\]
The $3x+1$ conjecture asserts that every orbit reaches $1$; see Lagarias
\cite{Lagarias1985} for the classical survey. Terras \cite{Terras1976}
introduced the \emph{stopping time}
\[
\sigma(n) \;=\; \inf\{\, t \ge 1 : T^{t}(n) < n \,\},
\]
proved that $\sigma(n)$ is finite for almost every $n$, and observed that the
iterate admits the exact affine form
\begin{equation}\label{eq:affine}
T^{t}(n) \;=\; \frac{3^{u_t(n)}\, n + b_t(n)}{2^{t}},
\end{equation}
where $u_t(n)$ counts the odd iterates among $n, T(n), \dots, T^{t-1}(n)$ and
$b_t(n) \ge 0$ depends only on the parity word of the first $t$ steps. The
\emph{coefficient stopping time} is
\[
\tau(n) \;=\; \inf\{\, t \ge 1 : 3^{u_t(n)} < 2^{t} \,\},
\]
the first time the multiplicative coefficient in \eqref{eq:affine} drops below
$1$. Since $b_t \ge 0$, a drop of the value forces a drop of the coefficient:
$\tau(n) \le \sigma(n)$ for all $n$. Terras conjectured that the combinatorial
quantity captures the arithmetic one exactly:

\begin{conjecture}[Coefficient stopping time conjecture, CST; Terras 1976]
\label{conj:cst}
$\tau(n) = \sigma(n)$ for every integer $n > 1$.
\end{conjecture}

The conjecture matters because $\tau$ is a function of the residue of $n$
modulo $2^{\tau}$ alone. If CST holds, the level sets
$\{n : \sigma(n) = t\}$ are exactly finite unions of arithmetic progressions
modulo $2^{t}$, and all stopping-time densities are exactly computable ---
this was Terras' motivation. The exclusion $n>1$ is forced: $\tau(1)=2$ while
$\sigma(1)=\infty$, because $1 \to 2 \to 1$ is the trivial cycle. As we will
see, this exception is not an accident but the exact \emph{equality case} of
the violation inequality, and the entire conjecture turns out to be governed
by how well $\lb = 1.5849625\ldots$ can be approximated from above by
rationals $j/u$ --- that is, by the continued fraction convergents of $\lb$,
whose first nontrivial upper member is the Pythagorean interval
$2^{8}/3^{5} = 256/243$.\footnote{Strictly, $256/243$ is the Pythagorean
\emph{limma} (the diatonic semitone); the comma proper is
$3^{12}/2^{19} = 531441/524288$, attached to the next convergent $19/12$, and
Mercator's comma $3^{53}/2^{84}$ belongs to $84/53$. We follow loose usage and
call all these small intervals ``commas.'' The from-below members ($19/12$,
$84/53$, \dots) satisfy $3^{u} > 2^{j}$ and are harmless for CST; the danger
zones are exactly the from-above members $2/1$, $8/5$, $65/41$, $485/306$,
$24727/15601$, \dots}

\subsection*{Results}

\begin{enumerate}
\item[(a)] CST holds for every $n>1$ with $\tau(n) \le 24$, by exhaustive
symbolic enumeration of all $81{,}119$ stopping classes of depth $\le 24$
(Proposition~\ref{prop:enum}). The only class whose minimal member does not
clear its violation threshold is the class of $n=1$, where the threshold is
attained with equality --- the trivial cycle.
\item[(b)] CST holds for every odd $1 < n \le 10^{6}$ by direct computation
(zero violations; maximal $\tau$ observed $176$;
Proposition~\ref{prop:direct}).
\item[(c)] We prove the unconditional bound $b/2^{j} < u/3$ for every stopping
class (Lemma~\ref{lem:bbound}). Combined with (b) this gives: CST holds for
\emph{all} $n > 1$ with $\tau(n) \le 4700$ (Corollary~\ref{cor:4700}), and in
every deeper stopping class except those whose $(u,j)$ lies in an explicit
sparse set of upper convergents and semiconvergents of $\lb$
(Theorem~\ref{thm:reduction}).
\item[(d)] The smallest safety margins concentrate at $(u,j) = (5,8)$, ratio
$256/243$; the closest observed approach to a violation is $n = 39$, which
after one comma word of eight steps lands on $38$, one unit below itself
(Section~\ref{sec:comma}).
\item[(e)] The extremal constant $B(t) = \max b/2^{j}$ over stopping classes
of depth $t$ is \emph{measured} to grow linearly, approximately one unit per
comma block ($1.25, 2.27, 3.24$ at $t = 8, 16, 24$), with extremal classes at
the repeated comma budgets $(5,8), (10,16), (15,24)$; we give a heuristic
sketch of the mechanism (Section~\ref{sec:bgrowth}). Only the proven bound of
Lemma~\ref{lem:bbound} is used elsewhere in the paper.
\item[(f)] CST is thereby reduced to \emph{effective rational approximation of
$\lb$}: lower bounds for $2^{j} - 3^{u}$ at the upper convergents, i.e.\ the
territory of Baker's theorem and Rhin's effective measures --- the identical
wall behind the exclusion of nontrivial cycles (Steiner \cite{Steiner1977},
Simons--de Weger \cite{SimonsDeWeger2005}, Hercher \cite{Hercher2023}).
Section~\ref{sec:onewall} makes the correspondence exact: a cycle attains the
violation threshold with equality, so CST violations are precisely
\emph{sub-threshold near-cycles}.
\end{enumerate}

Throughout we label every statement as PROVED (full proof supplied or cited),
VERIFIED (exhaustive computation, exact integer/rational arithmetic), MEASURED
(finite-range computation exhibiting a pattern), or SKETCH (heuristic
argument). The one genuinely heuristic item, the sharp $B$-growth law (e), is
quarantined: no other result depends on it.

\section{Preliminaries}\label{sec:prelim}

Fix $n \ge 1$ and let $(\epsilon_0, \epsilon_1, \dots)$,
$\epsilon_i \in \{0,1\}$, be the parity word: $\epsilon_i = T^{i}(n) \bmod 2$.
Iterating $T$ gives \eqref{eq:affine} with $u_t = \epsilon_0 + \cdots +
\epsilon_{t-1}$ and the constant $b_t$ obeying
\begin{equation}\label{eq:brec}
b_0 = 0, \qquad
b_{t+1} \;=\;
\begin{cases}
 b_t, & \epsilon_t = 0,\\
 3 b_t + 2^{t}, & \epsilon_t = 1.
\end{cases}
\end{equation}
Terras' basic lemma \cite{Terras1976} states that the first $t$ parities
depend only on $n \bmod 2^{t}$, and that the map from residues mod $2^{t}$ to
parity words of length $t$ is a bijection. Consequently, for each parity word
$w$ of length $j$ there is a unique residue $r \bmod 2^{j}$ realizing it, and
the triple $(u, j, b) = (u_j, j, b_j)$ is a function of $w$ alone.

\begin{definition}[Stopping class]
A parity word $w$ of length $j$ is a \emph{stopping word} if
$3^{u_t} \ge 2^{t}$ for all $t < j$ and $3^{u_j} < 2^{j}$. The associated
\emph{stopping class} is the residue class $\{ n \equiv r \pmod{2^{j}} \}$,
recorded as the state $(r, u, j, b)$. Every $n \ge 1$ with finite $\tau$ lies
in exactly one stopping class, and $\tau(n) = j$ for all members.
\end{definition}

For $j \ge 2$ every stopping word begins with $\epsilon_0 = 1$ (otherwise
$3^{0} < 2^{1}$ stops it at $t=1$), so all members of a stopping class of
depth $\ge 2$ are odd; conversely all even $n$ satisfy
$\tau(n)=\sigma(n)=1$. Hence CST is a statement about odd $n$.

\begin{lemma}[Violation inequality; PROVED]\label{lem:viol}
Let $(r,u,j,b)$ be a stopping class. Every member $n$ satisfies
$\sigma(n) \ge \tau(n) = j$, and
\[
\sigma(n) > \tau(n)
\quad\Longleftrightarrow\quad
n \;\le\; \frac{b}{2^{j} - 3^{u}} \;=:\; \Theta(w).
\]
\end{lemma}

\begin{proof}
For $t < j$ we have $3^{u_t} \ge 2^{t}$ and $b_t \ge 0$, so
$T^{t}(n) = (3^{u_t} n + b_t)/2^{t} \ge n$: no drop occurs before $j$, i.e.
$\sigma(n) \ge j$. At $t = j$,
$T^{j}(n) < n \iff 3^{u} n + b < 2^{j} n \iff n(2^{j}-3^{u}) > b$, and
$2^{j} > 3^{u}$ by definition of stopping.
\end{proof}

We call $\Theta(w)$ the \emph{violation threshold} of the class and
$r/\Theta(w)$ its \emph{margin}: CST holds in the class iff the margin of its
least member exceeds $1$ (then all members $r + k\,2^{j}$ exceed $\Theta$ as
well).

\begin{lemma}[The final step is a halving; PROVED]\label{lem:lastsGtep}
Every stopping word ends in $\epsilon_{j-1} = 0$, and consequently
\[
j \;=\; \ceil{\,u \lb\,}.
\]
Thus a stopping class' depth is determined by its odd-step count, and
\[
1 - \frac{3^{u}}{2^{j}} \;=\; 1 - 2^{-d(u)}, \qquad
d(u) := \ceil{\,u\lb\,} - u\lb \in (0,1).
\]
\end{lemma}

\begin{proof}
Suppose $\epsilon_{j-1} = 1$. Then $u = u_{j-1} + 1$ and, since the word did
not stop at $j-1$, $3^{u-1} \ge 2^{j-1}$, whence
$3^{u} \ge 3 \cdot 2^{j-1} > 2^{j}$, contradicting stopping. So
$\epsilon_{j-1}=0$, $u = u_{j-1}$, and $2^{j-1} \le 3^{u} < 2^{j}$, i.e.
$j - 1 \le u \lb < j$; since $\lb$ is irrational, $j = \ceil{u\lb}$.
\end{proof}

Lemma~\ref{lem:lastsGtep} is the hinge of the whole paper: the denominator
$2^{j} - 3^{u}$ of the violation threshold is small exactly when $d(u)$ is
small, i.e.\ when $u\lb$ lies just below an integer --- when $j/u$ is a good
rational approximation of $\lb$ \emph{from above}. By the classical theory of
one-sided best approximations (see e.g.\ Khinchin \cite{Khinchin}), such $u$
are precisely the denominators of the upper convergents and upper intermediate
fractions (semiconvergents) of the continued fraction
\[
\lb = [1; 1, 1, 2, 2, 3, 1, 5, 2, 23, 2, \dots],
\]
with convergents
$1,\ 2,\ \tfrac32,\ \tfrac85,\ \tfrac{19}{12},\ \tfrac{65}{41},\
\tfrac{84}{53},\ \tfrac{485}{306},\ \tfrac{1054}{665},\
\tfrac{24727}{15601},\ \dots$; the upper ones are
$\tfrac21, \tfrac85, \tfrac{65}{41}, \tfrac{485}{306},
\tfrac{24727}{15601}, \dots$

\section{Exhaustive verification}\label{sec:verify}

\subsection{Symbolic enumeration of stopping classes}

The recursion \eqref{eq:brec} together with Terras' bijection yields a direct
breadth-first enumeration. Maintain the set of \emph{survivor} states
$(r, u, j, b)$ (words with $3^{u_t} \ge 2^{t}$ for all $t \le j$), starting
from $(1, 0, 0, 0)$. Each survivor has two children $r' \in
\{r,\, r + 2^{j}\}$ modulo $2^{j+1}$; the child's next parity is the parity of
$(3^{u} r' + b)/2^{j}$, computable by modular arithmetic. A child with
$3^{u'} < 2^{j'}$ is a stopping class: record it, compute
$\Theta = b'/(2^{j'} - 3^{u'})$ in exact rational arithmetic, and if
$r' \le \Theta$ check the candidate $n = r'$ by direct iteration. All other
members of the class exceed $r'$ by multiples of $2^{j'} > \Theta$ in the
range explored, so $r'$ is the only candidate.

\begin{proposition}[VERIFIED]\label{prop:enum}
There are exactly $81{,}119$ stopping classes of depth $j \le 24$. In exactly
one of them the least member fails to exceed the violation threshold: the
class $(r,u,j,b) = (1, 1, 2, 1)$ of $n = 1$, where
$\Theta = 1/(2^{2}-3^{1}) = 1$ and the threshold is attained with
\emph{equality} --- the trivial cycle $1 \to 2 \to 1$. Consequently
\[
\tau(n) = \sigma(n) \quad \text{for every } n > 1 \text{ with } \tau(n)\le 24.
\]
The smallest margin $r/\Theta$ among all other classes is $2.020$.
\end{proposition}

(Computation: \texttt{scripts/70\_cst\_attack.py}; integer arithmetic
throughout, thresholds as exact fractions; re-run and confirmed for this
paper.)

\subsection{Direct verification to $10^{6}$}

\begin{proposition}[VERIFIED]\label{prop:direct}
For every odd $n$ with $1 < n \le 10^{6}$: $\tau(n) = \sigma(n)$. The maximal
coefficient stopping time observed is $\tau = 176$. Together with the trivial
even case, CST holds for all $1 < n \le 10^{6}$.
\end{proposition}

The two propositions overlap but are logically independent:
Proposition~\ref{prop:enum} covers all integers of small $\tau$ regardless of
size; Proposition~\ref{prop:direct} covers all small integers regardless of
$\tau$, and it is the second form --- a \emph{basis} --- that powers the
reduction of Section~\ref{sec:reduction}, because violation candidates are
forced to be small.

\section{The margins and the comma}\label{sec:comma}

Proposition~\ref{prop:enum} yields margin statistics over all $81{,}119$
classes. The ten smallest margins are:

\begin{center}
\begin{tabular}{crrr}
\hline
margin $r/\Theta$ & $(u,j)$ & $r$ & $2^{j}/3^{u}$ \\
\hline
$2.020$ & $(5,8)$   & $39$  & $1.05350$ \\
$2.391$ & $(3,5)$   & $11$  & $1.18519$ \\
$3.805$ & $(17,27)$ & $251$ & $1.03924$ \\
$3.965$ & $(5,8)$   & $79$  & $1.05350$ \\
$4.200$ & $(2,4)$   & $3$   & $1.77778$ \\
$4.507$ & $(4,7)$   & $7$   & $1.58025$ \\
$5.013$ & $(5,8)$   & $123$ & $1.05350$ \\
$5.853$ & $(5,8)$   & $95$  & $1.05350$ \\
$6.053$ & $(3,5)$   & $23$  & $1.18519$ \\
$9.141$ & $(5,8)$   & $199$ & $1.05350$ \\
\hline
\end{tabular}
\end{center}

Five of the ten smallest margins, including the minimum, live at
$(u,j) = (5,8)$: the Pythagorean interval
\[
\frac{2^{8}}{3^{5}} \;=\; \frac{256}{243} \;=\; 1.053497\ldots,
\]
the first nontrivial upper convergent $8/5$ of $\lb$. The minimal-margin
witness is concrete and pleasing: $n = 39$ has parity word $11101100$
($u=5$, $j=8$, $b=251$, $\Theta = 251/13 = 19.31$), and
\[
T^{8}(39) \;=\; \frac{3^{5}\cdot 39 + 251}{2^{8}} \;=\; \frac{9728}{256}
\;=\; 38.
\]
After one full comma word, $39$ lands one unit below itself. This is the
closest approach to a CST violation anywhere in the $81{,}119$ classes.

The mechanism is transparent from Lemmas \ref{lem:viol} and
\ref{lem:lastsGtep}: the threshold
$\Theta = (b/2^{j}) / (1 - 3^{u}/2^{j})$ explodes exactly when
$1 - 2^{-d(u)}$ is small, i.e.\ at upper convergents of $\lb$. The relevant
data for the first few are (exact rational arithmetic; threshold bound from
Lemma~\ref{lem:bbound} below):

\begin{center}
\begin{tabular}{crrrr}
\hline
$(u,j)$ & $2^{j}/3^{u}$ & $1 - 3^{u}/2^{j}$ & $\dfrac{u/3}{1-3^{u}/2^{j}}$ &
status \\
\hline
$(1,2)$        & $1.33333$   & $2.50\cdot 10^{-1}$ & $1.34$ &
 equality case $n=1$ \\
$(5,8)$        & $1.05350$   & $5.08\cdot 10^{-2}$ & $32.9$ &
 covered ($\ge 2.02\times$) \\
$(41,65)$      & $1.011529$  & $1.14\cdot 10^{-2}$ & $1.20\cdot 10^{3}$ &
 covered by basis \\
$(306,485)$    & $1.0010228$ & $1.02\cdot 10^{-3}$ & $9.99\cdot 10^{4}$ &
 covered by basis \\
$(15601,24727)$& $1.0000182$ & $1.82\cdot 10^{-5}$ & $2.86\cdot 10^{8}$ &
 open (needs larger basis) \\
\hline
\end{tabular}
\end{center}

The next danger zone after the comma is thus $(u,j) = (41,65)$ with
$2^{65}/3^{41} = 1.011529\ldots$; its threshold bound $1.2\cdot 10^{3}$ is
comfortably inside the verified basis $10^{6}$, as is that of
$(306,485)$.\footnote{An earlier internal note quoted $2^{65}/3^{41} \approx
1.0035$; the correct value, $36893488147419103232 / 36472996377170786403 =
1.011529\ldots$, is used here.} The first zone that escapes the $10^{6}$
basis under the proven bound is the semiconvergent $u = 2966$ (see
Corollary~\ref{cor:4700}), and the first \emph{convergent} to escape is
$(15601, 24727)$.

\section{The growth of $B(t)$}\label{sec:bgrowth}

The threshold $\Theta = (b/2^{j})/(1 - 3^{u}/2^{j})$ has two factors: the
Diophantine denominator, discussed above, and the numerator $b/2^{j}$. Define
\[
B(t) \;=\; \max\{\, b/2^{j} \;:\; (r,u,j,b) \text{ a stopping class with }
j = t \,\}.
\]

\begin{lemma}[Linear bound; PROVED]\label{lem:bbound}
For every stopping class $(r,u,j,b)$,
\[
\frac{b}{2^{j}} \;<\; \frac{u}{3},
\qquad\text{hence}\qquad
\Theta \;<\; \frac{u/3}{1 - 3^{u}2^{-j}}
\quad\text{and}\quad
B(t) \;<\; \frac{t}{3\lb} \;<\; 0.2104\, t .
\]
\end{lemma}

\begin{proof}
Set $y_t = b_t/3^{u_t}$, so $y_0 = 0$. By \eqref{eq:brec}, an even step leaves
$y$ unchanged and an odd step at depth $t$ gives
\[
y_{t+1} \;=\; \frac{3b_t + 2^{t}}{3^{u_t+1}} \;=\; y_t +
\frac{1}{3}\cdot\frac{2^{t}}{3^{u_t}} \;\le\; y_t + \frac13,
\]
since the survivor condition at depth $t$ ($3^{u_t} \ge 2^{t}$) holds for
every proper prefix of a stopping word. With $u$ odd steps in total,
$y_j \le u/3$. At the stopping depth, $3^{u} < 2^{j}$, so
\[
\frac{b}{2^{j}} \;=\; y_j \cdot \frac{3^{u}}{2^{j}} \;<\; y_j \;\le\;
\frac{u}{3}.
\]
For the global form, stopping gives $3^{u} < 2^{j}$, i.e.\
$u < j/\lb$.
\end{proof}

This settles, with room to spare, the bounding step that the reduction below
requires. The measured truth is sharper and structurally striking:

\begin{observation}[MEASURED, $t \le 27$]\label{obs:bgrowth}
$B(8) = 1.246$, $B(16) = 2.271$, $B(24) = 3.243$, $B(27) = 4.086$, attained at
the classes $(u,j) = (5,8), (10,16), (15,24), (17,27)$ respectively: the
extremal classes carry exactly the repeated comma budget $u/j \approx 5/8$,
and $B$ grows by approximately one unit per comma block,
$B(t) \approx 1 + t/8$.
\end{observation}

\begin{remark}[SKETCH of the mechanism]
In the proof of Lemma~\ref{lem:bbound} the gain of $y$ at an odd step of
depth $t$ is exactly $\tfrac13 \rho_t^{-1}$ with
$\rho_t = 3^{u_t}/2^{t} \ge 1$, and the final value is
$b/2^{j} = y_j\rho_j$ with $\rho_j < 1$. Both factors therefore reward
trajectories that hug criticality: gains are maximal ($\to \tfrac13$) when
$\rho \approx 1$, and the final discount $\rho_j$ is mildest when the word
stops barely below $1$. A survivor word that stays near-critical for $t$
steps must spend its odd steps at rate $\approx 1/\lb = 0.6309$, and the
integer word lengths realizing this most efficiently are the comma budgets
$(5k, 8k)$ (with the odd steps front-loaded so that every prefix survives).
Each 8-step comma block then contributes about
$5 \cdot \tfrac13 \cdot \langle \rho^{-1} \rangle \approx 1$ to $y$ while
$\rho$ decays only by the comma factor $3^{5}/2^{8} = 0.9494$, giving the
linear law with slope $\approx 1/8$. We emphasize: this paragraph is a
heuristic; the results of Section~\ref{sec:reduction} use only
Lemma~\ref{lem:bbound}.
\end{remark}

\section{The reduction theorem}\label{sec:reduction}

\begin{theorem}[Reduction to the upper convergents; PROVED, with
computational input]\label{thm:reduction}
Suppose $\tau(n) = \sigma(n)$ has been verified for all $1 < n \le N$
(as in Proposition~\ref{prop:direct} with $N = 10^{6}$). Then
$\tau(n) = \sigma(n)$ for every $n > 1$ lying in any stopping class
$(u, j)$, $j = \ceil{u\lb}$, with
\[
\frac{u/3}{1 - 2^{-d(u)}} \;<\; N, \qquad d(u) = \ceil{u\lb} - u\lb .
\]
In particular CST can fail only in classes whose $(u,j)$ makes $j/u$ a best
one-sided rational approximation of $\lb$ from above --- an upper convergent
or upper semiconvergent --- with $d(u)$ of order $u/N$ or smaller.
\end{theorem}

\begin{proof}
Let $n > 1$ violate CST in such a class. By Lemma~\ref{lem:viol} and
Lemma~\ref{lem:bbound},
\[
n \;\le\; \Theta \;<\; \frac{u/3}{1 - 3^{u}2^{-j}} \;=\;
\frac{u/3}{1 - 2^{-d(u)}} \;<\; N,
\]
using Lemma~\ref{lem:lastsGtep} for $j = \ceil{u\lb}$. Classes of depth
$j \ge 2$ consist of odd numbers, so $n$ is an odd integer in $(1, N]$,
contradicting the verified basis. For the final assertion: if
$\tfrac{u/3}{1-2^{-d(u)}} \ge N$ then $1 - 2^{-d(u)} \le u/(3N)$, so
$d(u) \le \log_2\bigl(1 - u/(3N)\bigr)^{-1} = O(u/N)$; small $d(u)$ forces
$j/u$ to approximate $\lb$ from above, and the one-sided best approximations
are the upper convergents and semiconvergents \cite{Khinchin}.
\end{proof}

\begin{corollary}[VERIFIED, exact rational arithmetic]\label{cor:4700}
With the basis $N = 10^{6}$:
\[
\frac{u/3}{1 - 2^{-d(u)}} < 10^{6} \quad\text{for every } u \le 2965,
\]
hence $\tau(n) = \sigma(n)$ for every $n > 1$ with $\tau(n) \le 4700$
(depths $j \le 4700$ correspond exactly to $u \le 2965$). The first
exceptional odd-step count is $u = 2966$ (a semiconvergent:
$2966 = 306 + 4\cdot 665$), with $j = 4701$ and threshold bound
$1.17 \cdot 10^{6}$, marginally beyond the basis. Among all
$u \le 40{,}000$ exactly $385$ values escape the $10^{6}$ basis, all of them
upper semiconvergent or convergent denominators; the largest threshold bound
in that range is $2.86\cdot 10^{8}$, at $u = 15601$ and $u = 2\cdot 15601$.
\end{corollary}

Three remarks on what the reduction does and does not achieve.

\emph{(i) The verified frontier moves from $\tau \le 24$ to $\tau \le 4700$.}
The enumeration of Proposition~\ref{prop:enum} costs exponential work in the
depth; the reduction replaces it, up to depth $4700$, by a one-dimensional
Diophantine scan plus the fixed basis. Extending the basis to
$3\cdot 10^{8}$ --- a modest computation by current standards (full Collatz
convergence has been machine-verified to $2^{68}$ \cite{Barina2021}, and a
$\tau$-$\sigma$ check stops at the first drop, on average after $\approx 3.5$
steps) --- would cover every $u \le 40{,}000$, i.e.\ every $n>1$ with
$\tau(n) \le 63{,}397$.

\emph{(ii) No finite basis finishes the job.} Along the upper convergents
$d(u) \to 0$ faster than any fixed multiple of $1/u$ infinitely often
(indeed $d(u_k) \asymp 1/(a_{k+1} u_k)$), so the thresholds
$\tfrac{u/3}{1-2^{-d(u)}} \asymp u^{2} a_{k+1}$ are unbounded: each new upper
convergent eventually demands a larger basis. The sequence of demands is,
however, only polynomial in the depth \emph{if} $\lb$ admits an effective
irrationality measure, which it does: Baker's theorem \cite{Baker1966} gives
effective lower bounds for $|j - u\lb|$, and Rhin's Pade-based measure
\cite{Rhin1987} gives, in the form used by Simons and de Weger
\cite{SimonsDeWeger2005}, a bound of the shape $d(u) > c\, u^{-\kappa}$ with
explicit $c$ and $\kappa$ ($\kappa = 13.3$ in their application). Through
Theorem~\ref{thm:reduction} this makes the basis requirement for depth $T$
effectively computable and polynomial in $T$, roughly $O(T^{\kappa + 1})$:
CST up to any prescribed depth is a \emph{finite, effectively bounded}
computation. What no current bound delivers is a uniform closure: full CST
would follow from $d(u) \ge c\,u/N_0$ for a fixed verified basis $N_0$ and
all large $u$ --- a strength of effective approximation far beyond present
technology.

\emph{(iii) The failure set is thin and explicit.} Any counterexample to CST
must be an odd $n$, larger than every verified basis, sitting inside a
specific residue class modulo $2^{j}$ whose modulus $2^{j}$ is
\emph{astronomically} larger than the violation threshold $O(u^{2}a)$: the
class contains at most one candidate below its threshold, namely its least
element $r$, and a violation requires the arithmetically determined $r$ to be
freakishly small ($r < \Theta$ while $r$ ranges over $[1, 2^{j}]$;
heuristically probability $\Theta/2^{j} \approx 0$). CST failures are exactly
\emph{near-cycles}, as the next section makes precise.

\section{One wall, two sides: CST and cycle exclusion}\label{sec:onewall}

Suppose $n$ is a periodic point of $T$ with period $j$ and $u$ odd steps per
period. Then \eqref{eq:affine} gives $n = (3^{u} n + b)/2^{j}$, i.e.
\begin{equation}\label{eq:cycle}
n \;=\; \frac{b}{2^{j} - 3^{u}} \;=\; \Theta(w):
\end{equation}
\emph{a positive cycle attains the violation threshold of its parity word
with equality}. The trivial cycle is the instance $(u,j,b) = (1,2,1)$,
$\Theta = 1$ --- precisely the unique sub-threshold member found in
Proposition~\ref{prop:enum}. Lemma~\ref{lem:viol} says that a CST violation
is the inequality relaxation of \eqref{eq:cycle}: an integer that does not
close the loop but sits at or below the rational fixed point of its own
word.\footnote{For a violating $n$, $T^{j}(n) \ge n$ with the coefficient
already below $1$: the orbit has failed to fall despite a falling word ---
it is shadowing the rational fixed point $\Theta$ from below.}

The known partial results on cycles run through lower bounds on the same
denominator $2^{j} - 3^{u}$. Steiner \cite{Steiner1977} excluded nontrivial
$1$-circuit cycles using Baker's theorem; Simons and de Weger
\cite{SimonsDeWeger2005} excluded $m$-circuit cycles for $m \le 68$ using
Rhin's effective measure \cite{Rhin1987}; Hercher \cite{Hercher2023} pushed
this to $m \le 91$. In every case the engine is: a cycle forces
$(2^{j}-3^{u}) \mid b$ with $b$ of size comparable to $2^{j}$, which is
incompatible with effective lower bounds for $|2^{j}-3^{u}|$ unless
$(u,j)$ is enormous. Theorem~\ref{thm:reduction} shows that CST consumes the
identical input: effective lower bounds for $2^{j} - 3^{u}$ at upper
convergents of $\lb$. The two classical subproblems of the $3x+1$ problem
that admit partial rigorous progress --- ``no nontrivial cycles'' and
``stopping is decided by the parity word'' (CST) --- are thus a single
question about the effective rational approximation of $\lb$, seen from two
sides: equality (divisibility) for cycles, inequality (smallness of the
least class member) for CST. The Pythagorean comma $256/243$ is the first
nontrivial post of that wall; the margin $2.02$ of Section~\ref{sec:comma}
is how close the integers come to it at small depth.

\section{Conclusion and honest status}\label{sec:conclusion}

We have verified the Terras coefficient stopping conjecture exhaustively in
depth ($\tau \le 24$, all $81{,}119$ stopping classes) and in height (all
$n \le 10^{6}$), proved the linear bound $B(t) < t/(3\lb)$ on the additive
constants, and combined the two into an unconditional extension of the
verified frontier to $\tau \le 4700$ together with a complete localization of
any possible failure: CST can fail only at explicitly listed upper
convergents and semiconvergents of $\lb$, only at integers below an
explicitly computable threshold $O(u^{2} a_{k+1})$, and only in the guise of
a sub-threshold near-cycle. The conjecture is thereby reduced to effective
rational approximation of $\lb$ --- Baker--Rhin territory --- the same wall
that bounds cycle exclusion.

Status of each component:

\begin{center}
\begin{tabular}{lll}
\hline
Statement & Location & Status \\
\hline
$\tau \le \sigma$; violation inequality & Lemma~\ref{lem:viol} & proved \\
Final halving; $j = \ceil{u\lb}$ & Lemma~\ref{lem:lastsGtep} & proved \\
CST for $\tau(n) \le 24$ & Prop.~\ref{prop:enum} &
 verified (exhaustive, exact) \\
CST for $1 < n \le 10^{6}$ & Prop.~\ref{prop:direct} &
 verified (exhaustive) \\
Margin concentration at $256/243$ & Sec.~\ref{sec:comma} &
 verified (finite range) \\
$B(t) < t/(3\lb)$ & Lemma~\ref{lem:bbound} & proved \\
$B(t) \approx 1 + t/8$, comma extremizers & Obs.~\ref{obs:bgrowth} &
 measured ($t \le 27$) + sketch \\
Reduction theorem & Thm.~\ref{thm:reduction} & proved (uses
 Prop.~\ref{prop:direct}) \\
CST for $\tau(n) \le 4700$ & Cor.~\ref{cor:4700} &
 proved + verified scan \\
Stitching via Rhin-type measures & Sec.~\ref{sec:reduction}(ii) &
 template, not carried out \\
\hline
\end{tabular}
\end{center}

Two computations would extend the frontier mechanically: (a) a basis to
$3\cdot 10^{8}$ lifts Corollary~\ref{cor:4700} to $\tau \le 63{,}397$; (b)
instantiating Rhin's explicit constants in Theorem~\ref{thm:reduction} turns
``CST up to depth $T$'' into an explicit function of feasible computation for
every $T$. The step that would close CST outright --- a uniform effective
bound $d(u) \gg u/N_0$ --- is of the same order of difficulty as excluding
all cycles, and we do not claim it. What this note establishes is that both
problems now wait at the same door, and that the door is labeled
$256/243$.

\subsection*{Computational note}

All enumerations use exact integer and rational arithmetic (Python,
\texttt{fractions.Fraction}); the stopping-class enumeration is
\texttt{scripts/70\_cst\_attack.py} in the authors' repository, and every
number quoted in this paper (class counts, margins, $B(t)$ values, convergent
gaps, the boundary cases $u = 2965, 2966$ of Corollary~\ref{cor:4700}) was
recomputed independently for the final text.

\begin{thebibliography}{9}

\bibitem{Terras1976}
R.~Terras,
\emph{A stopping time problem on the positive integers},
Acta Arithmetica \textbf{30} (1976), 241--252.

\bibitem{Lagarias1985}
J.~C.~Lagarias,
\emph{The $3x+1$ problem and its generalizations},
American Mathematical Monthly \textbf{92} (1985), 3--23.

\bibitem{Steiner1977}
R.~P.~Steiner,
\emph{A theorem on the Syracuse problem},
Proceedings of the 7th Manitoba Conference on Numerical Mathematics and
Computing (Winnipeg, 1977), 553--559.

\bibitem{SimonsDeWeger2005}
J.~Simons and B.~de~Weger,
\emph{Theoretical and computational bounds for $m$-cycles of the $3n+1$
problem},
Acta Arithmetica \textbf{117} (2005), 51--70.

\bibitem{Hercher2023}
C.~Hercher,
\emph{There are no Collatz $m$-cycles with $m \le 91$},
arXiv:2201.00406 (2023).

\bibitem{Baker1966}
A.~Baker,
\emph{Linear forms in the logarithms of algebraic numbers I},
Mathematika \textbf{13} (1966), 204--216.

\bibitem{Rhin1987}
G.~Rhin,
\emph{Approximants de Pad\'e et mesures effectives d'irrationalit\'e},
S\'eminaire de Th\'eorie des Nombres, Paris 1985--86,
Progress in Mathematics \textbf{71}, Birkh\"auser (1987), 155--164.

\bibitem{Khinchin}
A.~Ya.~Khinchin,
\emph{Continued Fractions},
University of Chicago Press, 1964.

\bibitem{Barina2021}
D.~Barina,
\emph{Convergence verification of the Collatz problem},
The Journal of Supercomputing \textbf{77} (2021), 2681--2688.

\end{thebibliography}

\end{document}
